\scanpage{84}
\mspage{6}

\deleted{\unsure{On définit de|On déduit de|On définit des}}

Le \deleted{catégorie} \inserted{champ} en question n'est pas rigide -- i.e. un objet peut avoir des \unsure{autom|automs} non-triviaux. \deleted{Les} \inserted{Un} autom \deleted{sont} \inserted{d'un objet (34) \unsure{est|en}} défini\deleted{\unsure{s}} par un autom de $\underline J$, i.e. une section du schéma en groupes étale $\unsure{\mathfrak S|\mathcal G}_{\underline J}$, qui soit compatible avec les données b) et c). La compatibilité avec c) signifie que c'est une section de

\[
\mathfrak S_{\underline J}^{+}\simeq\deleted{\unsure{\mathrm{id}|\mathbb Z}}
\underbrace{S}_{\substack{\text{correspond}\\\text{à la section}\\\text{unité du}\\\text{schéma en groupes}}}
\amalg\underline\omega_{\underline J}
\]

une section qui est $\ne\mathrm{id}$ sur tte fibre géom. revient donc à une section de $\underline\omega_{\underline J}$, soit $\rho$ -- de sorte que $\underline J$ devient un torseur sous $(\mathbb Z/3\mathbb Z)_{S\deleted{\unsure{3}}}$. La compatibilité avec b) signifie que $t=\rho(t)$ i.e. que $t$ \deleted{correspond à} \inserted{est une} section du sous-schéma

\diagram{\begin{tikzcd}[ampersand replacement=\&,row sep=small,column sep=small]
\underline P_{\underline J}\arrow[d,equal] \&[-1ex] {\subset} \&[-1ex] \Sigma_{\underline J}^*,\\
|[draw,ellipse]|(\Sigma_{\underline J})^{\mathfrak S_{\underline J}^{+}}
\end{tikzcd}}

\deleted{\unsure{partout|partant} dès que}
% archive sheet number: 83
